% ============================================================
%  APENDICES A-D (2026-11-24: un apéndice por tema, con títulos descriptivos)
%  Bloques movidos LITERALMENTE desde ch3_marco_teorico.tex el 2026-11-23
%  (observacion 3 del director). Solo cambian: niveles de encabezado,
%  etiquetas app:*, dos correcciones objetivas en A.4 (raiz sexta en
%  eq:tau_min; ~4000 transitos segun Keller Fig. 6.17) y la atribucion
%  inicial de A.4. Detalle: drafts/cap3/CAMBIOS.md
% ============================================================

\chapter{Propagación del pulso: dispersión y efecto Kerr}
\label{app:propagacion}

\section{La envolvente de campo lentamente variable}

Un pulso ultracorto es una oscilación óptica de frecuencia portadora $\omega_0$
modulada por una envolvente que varía mucho más despacio que la portadora. Esta
separación de escalas justifica la \emph{aproximación de envolvente lentamente
variable}, en la que el campo eléctrico real se escribe como
%
\begin{equation}
    E(z,t) = \tfrac{1}{2}\,A(z,t)\,
    e^{\,i(\omega_0 t - k_0 z)} + \text{c.c.},
    \label{eq:envolvente}
\end{equation}
%
donde $A(z,t)$ es la envolvente compleja, cuyo módulo $|A|^2$ mide la intensidad
instantánea y cuya fase codifica la estructura de frecuencias del pulso, $k_0$
es el número de onda a la portadora y c.c.\ denota el complejo conjugado. Toda la
dinámica de interés (ensanchamiento, \emph{chirp}, acortamiento) se traslada así a
la evolución de $A$, mucho más lenta que las oscilaciones a $\omega_0$. Describir
esa evolución es el objeto de esta sección.

\section{Régimen lineal: dispersión de la velocidad de grupo}

En un medio transparente, las distintas componentes espectrales del pulso viajan
con velocidades de fase diferentes porque el número de onda depende de la
frecuencia, $k=k(\omega)$. Desarrollando esta relación de dispersión en torno a la
portadora,
%
\begin{equation}
    k(\omega) = k_0 + k_1\,(\omega-\omega_0)
    + \tfrac{1}{2}\,k_2\,(\omega-\omega_0)^2 + \cdots,
    \qquad
    k_n \equiv \left.\frac{d^{\,n}k}{d\omega^{\,n}}\right|_{\omega_0},
    \label{eq:beta_expansion}
\end{equation}
%
se identifican dos coeficientes de significado físico directo. El término de
primer orden, $k_1 = 1/v_g$, es el inverso de la velocidad de grupo y no deforma
el pulso: se limita a desplazarlo en bloque. Adoptando un marco de referencia que
viaja con el pulso, $t \to t - z/v_g$, ese término se elimina y la deformación
queda gobernada por el coeficiente de segundo orden $k_2$, la \emph{dispersión de
la velocidad de grupo} (GVD). En este marco comóvil, la envolvente evoluciona
según
%
\begin{equation}
    \frac{\partial A}{\partial z}
    = i\,D_0\,\frac{\partial^2 A}{\partial t^2},
    \qquad D_0 \equiv \tfrac{1}{2}\,k_2,
    \label{eq:gvd_prop}
\end{equation}
%
una ecuación de tipo Schrödinger libre cuyo efecto es puramente dispersivo: no
altera el espectro $|\tilde A(\omega)|^2$, sino que reorganiza las fases
relativas de sus componentes. El signo de $k_2$ separa dos regímenes. En
\emph{dispersión normal} ($k_2>0$, es decir $D_0>0$) las componentes de baja
frecuencia adelantan a las de alta, y el pulso adquiere un \emph{chirp} positivo
al tiempo que se ensancha temporalmente. En \emph{dispersión anómala}
($k_2<0$, $D_0<0$) ocurre lo contrario. Este segundo régimen es el relevante para
la formación de solitones, como se verá en la Sección~\ref{sec:solitones}: solo
con dispersión anómala puede la no linealidad de Kerr contrarrestar el
ensanchamiento y estabilizar el pulso.

Conviene subrayar que la GVD, por sí sola, conserva el ancho de banda del pulso y
por tanto su duración mínima alcanzable. Un pulso cuya duración temporal es la
menor compatible con su ancho espectral se denomina \emph{limitado por
transformada}; la dispersión lo aleja de ese límite introduciendo \emph{chirp},
pero una compensación dispersiva de signo opuesto puede, en principio,
restaurarlo. Esta reversibilidad distingue a la dispersión de los procesos que sí
modifican el contenido espectral, a los que se pasa a continuación.

\section{Régimen no lineal: efecto Kerr y automodulación de fase}

A las intensidades pico que alcanzan los pulsos ultracortos, la respuesta del
medio deja de ser lineal. El efecto Kerr óptico hace que el índice de refracción
dependa de la intensidad instantánea,
%
\begin{equation}
    n(I) = n_0 + n_2\,I,
    \qquad I \propto |A|^2,
    \label{eq:kerr}
\end{equation}
%
con $n_2>0$ el índice no lineal del material \cite[sec.~4.1.1]{Keller2022}. Puesto que la fase
acumulada por el campo es proporcional al índice, una envolvente de intensidad
variable imprime sobre sí misma un desfase que sigue su propio perfil temporal:
este es el fenómeno de \emph{automodulación de fase} (SPM). En términos de la
envolvente, su contribución a la propagación es
%
\begin{equation}
    \frac{\partial A}{\partial z}
    = -\,i\,\delta_0\,|A|^2 A,
    \qquad \delta_0 \propto n_2\,\omega_0,
    \label{eq:spm_prop}
\end{equation}
%
donde $\delta_0>0$ es el coeficiente de no linealidad. A diferencia de la GVD, la
SPM no reordena fases a espectro constante: al ser el desfase proporcional a la
intensidad, sus flancos temporales, donde $|A|^2$ varía con rapidez, generan
\emph{nuevas} componentes de frecuencia, ensanchando el espectro del pulso. Es,
por tanto, el mecanismo que permite acceder a anchos de banda, y con ellos a
duraciones, que la sola cavidad lineal no alcanzaría.

\chapter{Promediado por ronda de la ecuación maestra de Haus}
\label{app:promediado}

El punto de partida es la separación de dos escalas temporales. Por un lado, el
\emph{tiempo local} $t$, definido en un marco de referencia que viaja con el
pulso, resuelve la estructura interna de la envolvente (la escala de
femto\-/picosegundos). Por otro, el \emph{tiempo lento} $T$ mide la evolución de
esa envolvente a lo largo de muchos tránsitos y avanza en pasos del tiempo de ida
y vuelta $T_R$; resulta cómodo tratarlo como una variable continua, de modo que
$T = n\,T_R$ con $n$ el número de ronda. La cantidad de interés es entonces la
envolvente compleja $A(T,t)$, lentamente variable en $T$ y rápidamente
estructurada en $t$.

Cada elemento intracavitario $j$ produce, en un solo paso, un cambio
$\Delta_j A$ sobre la envolvente. Tras un recorrido completo, el campo acumula la
suma de todas esas contribuciones:
%
\begin{equation}
    A(T + T_R,\, t) = A(T,t) + \sum_j \Delta_j A .
    \label{eq:haus_acumulacion}
\end{equation}
%
En el régimen estacionario de bloqueo de modos, cada $\Delta_j A$ es pequeño
frente a $A$, y también lo es su suma. Bajo esta hipótesis de cambio débil por
ronda, el primer miembro de la ecuación~\eqref{eq:haus_acumulacion} admite el
desarrollo $A(T+T_R,t) \approx A(T,t) + T_R\,\partial_T A$, con lo que la
diferencia entre dos vueltas consecutivas se convierte en una derivada:
%
\begin{equation}
    T_R\,\frac{\partial A}{\partial T} = \sum_j \Delta_j A .
    \label{eq:haus_maestra_forma}
\end{equation}
%
La ecuación~\eqref{eq:haus_maestra_forma} es la forma genérica de la ecuación
maestra: traduce la dinámica discreta «vuelta a vuelta» en una evolución
continua gobernada por la acción promediada de todos los elementos de la cavidad.
Los operadores $\Delta_j A$ de cada elemento se especifican en la
Sección~\ref{sec:haus}, a partir de los procesos físicos resumidos en el
Cuadro~\ref{tab:procesos}.

\chapter{Solución de Kuizenga--Siegman del bloqueo de modos activo}
\label{app:ks}

El punto de partida es la ecuación maestra para el bloqueo de modos activo sin
dispersión ni no linealidad, deducida en la Sección~\ref{sec:haus}
(ecuación~\eqref{eq:haus_activo_base}). Puesto que el pulso es mucho más corto que
el periodo de modulación $2\pi/\omega_m$, basta desarrollar el modulador en torno
a su mínimo de pérdida, $M\bigl(1-\cos\omega_m t\bigr)\approx M_s\,t^2$ con
$M_s\equiv M\omega_m^2/2$. En estado estacionario ($\partial_T A = 0$) la ecuación
maestra se reduce al problema de autovalores
%
\begin{equation}
    D_g\,\frac{\partial^2 A}{\partial t^2}
    - M_s\,t^2\,A + (g-l)\,A = 0,
    \label{eq:ks_autovalores}
\end{equation}
%
formalmente idéntico al del oscilador armónico cuántico. Su modo fundamental es un
pulso de perfil gaussiano,
%
\begin{equation}
    A(t) = A_0\,\exp\!\left(-\frac{t^2}{2\tau_g^{2}}\right),
    \qquad
    \tau_g = \left(\frac{D_g}{M_s}\right)^{1/4},
    \label{eq:ks_gaussiano}
\end{equation}
%
cuya anchura resulta del equilibrio entre los dos términos que compiten en
\eqref{eq:ks_autovalores}: la dispersión de ganancia $D_g=g/\Omega_g^2$, que
ensancha el pulso al filtrar sus alas espectrales, y la curvatura del modulador
$M_s$, que lo estrecha al penalizar sus alas temporales
(Figura~\ref{fig:ks-narrowing}).

 Sustituyendo
\eqref{eq:ks_gaussiano} en \eqref{eq:ks_autovalores} se obtiene además la
condición de balance que fija la ganancia saturada,
%
\begin{equation}
    g - l = \sqrt{D_g\,M_s},
    \label{eq:ks_balance}
\end{equation}
%
es decir, el exceso de ganancia sobre las pérdidas lineales necesario para
sostener el pulso estacionario.

\chapter{Teoría perturbativa del solitón en régimen activo}
\label{app:soliton_activo}

% [REV-2026-11-23] atribución explícita (antes solo había cita al final del pasaje)
Esta deducción sigue la exposición de Keller~\cite[sec.~6.6]{Keller2022}, de la que se
toman el planteamiento, las ecuaciones y los valores numéricos.

Se propone así un \emph{ansatz} que describe un solitón acompañado de un fondo
residual (el continuo),
%
\begin{equation}
    A(T,t) = A_0\,\operatorname{sech}\!\left(\frac{t}{\tau}\right)
    \exp\!\left(i\phi_0\,\frac{T}{T_R}\right) + \text{continuo},
    \label{eq:ansatz_soliton}
\end{equation}
%
donde $A_0$ es la amplitud de la envolvente, $\tau$ su duración y $\phi_0$ el
desplazamiento de fase que el solitón acumula por cada periodo de repetición
$T_R$. Con dispersión anómala suficiente, el continuo decae exponencialmente y,
para una energía de pulso dada, se forma un solitón cuya duración ideal es
%
\begin{equation}
    \tau = \frac{4\,|D|}{F_{P,L}\,\delta},
    \label{eq:tau_soliton}
\end{equation}
%
con $\delta$ el coeficiente de SPM y $F_{P,L}$ la densidad de energía del pulso en
el medio láser, bajo el supuesto de que la SPM se genera en el cristal de
ganancia. La dispersión neta por ronda $D$ se relaciona con la dispersión de
retardo de grupo acumulada mediante $\mathrm{GDD}=2D$, de modo que el régimen
anómalo ($\mathrm{GDD}<0$) corresponde, coherentemente con la
Sección~\ref{sec:propagacion}, a $D<0$. La ecuación~\eqref{eq:tau_soliton} indica
que reducir la dispersión acorta el solitón; sin embargo, como se verá, existe un
límite por debajo del cual el pulso se desestabiliza.

\paragraph{El continuo y su papel estabilizador.}
El solitón no se forma de manera instantánea. Partiendo de un pulso gaussiano
inicial, el sistema requiere del orden de miles de tránsitos (unos 4000 en el caso
de Nd:YAG de \cite[Fig.~6.17]{Keller2022}) para relajar
hacia el estado estacionario solitónico, de duración mucho menor. La fracción de
energía del pulso inicial que no logra ajustarse al perfil del solitón se vierte
en el continuo: una componente de baja intensidad para la cual la SPM es
insuficiente para compensar el ensanchamiento dispersivo, de modo que se dispersa
temporalmente. Al ensancharse, el continuo pasa más tiempo fuera de la ventana de
transmisión del modulador y sufre en él una pérdida creciente, por lo que decae
con las rondas (Figura~\ref{fig:saturable_gain}). En el marco de referencia del
solitón, el efecto Kerr le imprime un desplazamiento de fase periódico que genera
interferencias con el continuo; estas se desvanecen gradualmente a medida que el
continuo se extingue, dejando únicamente el solitón.



\paragraph{Duración mínima y factor de reducción.}
Existe un límite a cuánto puede acortarse el solitón manteniéndolo estable frente
al continuo: si la dispersión es demasiado baja, el continuo no se ensancha con la
rapidez suficiente para ser absorbido por el modulador y el sistema se
desestabiliza, llegando a generar pulsos satélite tras muchos miles de tránsitos.
La teoría perturbativa de solitones fija ese umbral a través de la dispersión
normalizada mínima
%
\begin{equation}
    D_{n,\min} = \frac{|D|}{D_g}
    = \frac{2}{9}\sqrt[3]{\frac{(9\phi_0/2)^2}{D_g\,M_s}},
    \label{eq:disp_min}
\end{equation}
%
a la que corresponde una duración a media altura de la intensidad
%
\begin{equation}
    \tau_{\min} = 1.763\,\sqrt[6]{\frac{2\,D_g^{2}}{9\,\phi_0\,M_s}}.
    \label{eq:tau_min}
\end{equation}
%
Comparando la duración mínima solitónica~\eqref{eq:tau_min} con la del pulso
gaussiano puramente activo~\eqref{eq:tau_gauss}, el factor de reducción que aporta
el régimen solitónico es
%
\begin{equation}
    R_{\max} = \frac{\tau_p}{\tau_{\min}}
    = \frac{1.665}{1.763}\,\sqrt[12]{\frac{(9\phi_0/2)^2}{D_g\,M_s}}
    = \frac{1.665}{1.763}\,\sqrt[4]{\tfrac{9}{2}\,D_{n,\min}}.
    \label{eq:factor_reduccion}
\end{equation}
%
A diferencia del caso puramente activo, la envolvente estacionaria ya no es un
pulso gaussiano con \emph{chirp}, sino un solitón $\operatorname{sech}^2$ sin
\emph{chirp}, que es lo que permite descender al régimen de femtosegundos.

\paragraph{Balance de energía y estabilidad.}
El solitón permanece estable mientras su pérdida efectiva $l_s$ sea inferior a la
que sufre el continuo $l_c$. La pérdida del solitón consta de dos aportaciones,
%
\begin{equation}
    l_s = \frac{D_g}{3\,\tau^2}
        + \frac{\pi^2}{24}\,M\,\omega_m^2\,\tau^2,
    \label{eq:perdida_soliton}
\end{equation}
%
la primera debida al ancho de banda finito de la ganancia, que tiende a ensanchar
el solitón, y la segunda a la acción residual del modulador de pérdida. El balance
de ganancia y pérdidas se escribe entonces como $g = l + l_s$, con $l$ la pérdida
aportada por el resto de los elementos de la cavidad. Para solitones muy cortos, el
primer término domina y la influencia del modulador sobre el solitón resulta casi
nula; el continuo, en cambio, sí experimenta en el modulador una pérdida intensa,
lo que garantiza su supresión y, con ella, la estabilidad del pulso solitónico.
Este contraste entre la escasa pérdida del solitón y la elevada del continuo es,
en última instancia, el mecanismo que sostiene el bloqueo de modos activo en su
régimen solitónico \cite{Haus2000,Keller2022}.
